%% notes-tex/figure-3-gate/sections/11-inductive.tex
%% SOURCE: tables/morph_v25.tex and tables/wandb_runs_morph_v25.tex are generated by
%% experiments/019-simb-multimodal/scripts/morph_v25_readout.py from results/morph_v25_readout.json
%% (W&B zhao-group/torchcell_019_morph_v25, read 2026-10-13 after the last cabbi task of arrays
%% 2428203, 2428214, 2428289, 2428290, 2428291, 2428295, 2428524 and 2428590 and the Delta packs
%% 22842101 to 22842108 finished); tables/prot_v26.tex and tables/wandb_runs_prot_v26.tex by
%% prot_v26_readout.py from results/prot_v26_readout.json (W&B torchcell_019_prot_v26, cabbi array
%% 2428296). The sequence ridge bars are results/morphology_baselines_split_fig3_core_moving/
%% (expression_baselines_split.py, job 3920); the v24 numbers are tables/morph_v24.tex.

\clearpage
\section{2026-10-13: the inductive case, rounds v25 and v26\secstatus{todo}}
\label{sec:inductive}

\begin{keybox}
\textbf{From the genotype alone the model now predicts the 116 moving CalMorph features at
$0.149 \pm 0.022$} (mean validation Pearson per feature over epochs 400 to 499, twelve runs:
four split seeds, three initializations, \cref{tab:morph-v25}), against $0.046 \pm 0.015$ for the
v24 readout on the same strains and the same features, a sequence-embedding ridge bar of 0.08
to 0.14, and $0.145 \pm 0.030$ for the transductive arm that reads the strain's own proteome
(v24, \cref{tab:morph-v24}). The cell is \texttt{M\_zs} with the MSE loss at lr $3\times10^{-4}$: the
global head reads the perturbed CLS concatenated with $z_S$, the strain's own perturbed gene
token (the PerturbationHead input behind the trigenic result), trained with MSE in place of
the 19-quantile pinball.

\textbf{What the round was about.} Every v24 genotype-only run sat on a plateau: by epoch 50
the head predicts the per-feature marginal (training loss 0.54 to 0.20) and the training
Pearson stays at 0.04 to 0.12 for 1{,}200 epochs; two runs of 27 left it. v25 and v26 varied
the readout, the loss, the learning rate, the graph prior, the batch and the schedule, 41
cells and 183 runs, scored on the registered window and on the escape fraction (the share of
runs whose eval-mode training Pearson passed 0.3).

\textbf{Three levers break the plateau, and they do not add.} The gene-token readout with
MSE (12 of 12 runs, escapes at epochs 180 to 410); a learning rate of $10^{-4}$ on the plain CLS
(7 of 8, $0.131 \pm 0.047$); the graph KL regularizer in place of the hard attention mask on
the CLS (4 of 4, $0.093 \pm 0.036$). MSE on the CLS (0 of 8 at two rates) or on the
cross-attention readout (1 of 4) does not escape; the pinball on $z_S$ escapes (8 of 10) but
late and unevenly ($0.097 \pm 0.072$); every cell at lr $10^{-3}$ or above stayed on the
marginal (0 of 27 cells, 0.03 to 0.04); CRPS, warm restarts, warmup and the mean gene pool
changed nothing.

\textbf{The contrastive arm is null to negative, measured.} Pulling the cosine similarity of
two strains' latents toward the correlation of their label vectors (\cref{sec:align}) left
every morphology cell on the marginal (0 of 20 escapes, 0.039 to 0.047) and suppressed the
\texttt{M\_zs} escapes (0 of 8 against 4 of 4 without it); on the knockout proteome it reads
$-0.008$ and $-0.010$ paired against a reference of $0.062 \pm 0.010$ (\cref{tab:prot-v26}).
The arm is dropped.

\textbf{Budget.} The loss stops moving by epoch 50, an escape happens by epoch 400, and an
escaped run peaks between epochs 400 and 600 and overfits after: the lr $10^{-4}$ cells run to
1{,}000 epochs scored lower over epochs 900 to 999 than over 400 to 499 in 7 of 8 runs. 500
epochs and the 400 to 499 window are the registered rule for this round, declared on
2026-10-12 before any cell was read on it.
\end{keybox}

\pfstep{The design} Round v25 (\file{igb_expr_wave5.slurm}, stages \texttt{morph\_v25},
\texttt{morph\_v25b}, \texttt{morph\_v25c}; Delta \file{delta_morph_v25_small.slurm}) is a
grid on the genotype-only arm of v24: the readout (\texttt{M\_cls}, the perturbed CLS;
\texttt{M\_pool}, the CLS plus the mean gene pool; \texttt{M\_zs}, the CLS plus $z_S$;
\texttt{M\_xattn}, one learned query per feature over every token) at lr $3\times10^{-4}$,
$10^{-3}$ and $3\times10^{-3}$ on split seeds 0 to 3; initialization replicates of
\texttt{M\_cls}; the loss (pinball, MSE, CRPS, Pearson plus MSE) on \texttt{M\_cls} and
\texttt{M\_zs}; the graph KL at $10^{-2}$ and 1 in place of the hard mask; lr $10^{-4}$; warm
restarts; a stacked cell; and, on Delta, batch 32, 64 and 128 with initialization replicates of
\texttt{M\_pool} and \texttt{M\_zs}. Round v26 added the contrastive term and, as v26b, the
combinations the first read pointed at: initialization replicates of \texttt{M\_zs} with MSE,
MSE on \texttt{M\_xattn}, MSE with the alignment term, and the lr $10^{-4}$ cells at 1{,}000
epochs. One run per card throughout, batch 128, config \file{cgt_morph_v24} with overrides;
every run is listed by cell in \cref{tab:wandb-runs-morph-v25}, the round's view is
\href{https://wandb.ai/zhao-group/torchcell_019_morph_v25?nw=xbewyhsxbtr}{torchcell\_019\_morph\_v25}.

%% GENERATED by experiments/019-simb-multimodal/scripts/morph_v25_readout.py. Do not edit by hand.
%% SOURCE: experiments/019-simb-multimodal/results/morph_v25_readout.json
\begin{table}[htbp]
  \centering
  \footnotesize
  \caption[Round v25, inductive morphology]{Round v25 by cell (arm plus recipe), genotype only on fig3_core, 116 moving features. The window is the mean validation Pearson per feature over epochs 400 to 499 of finished runs (the 1{,}200-epoch runs of the first wave also carry their 1{,}000 to 1{,}199 mean in the json); the escape fraction is the share of runs whose training Pearson exceeded 0.3 at any epoch; the training fit is at the last logged epoch. A cell with unfinished runs is partial. Bold, the best finished window mean. \src{experiments/019-simb-multimodal/scripts/morph_v25_readout.py}}
  \label{tab:morph-v25}
  \begin{tabular}{@{}lrrrrr@{}}
    \toprule
    cell & runs (n) & finished (n) & window mean val Pearson per feature (s.d.) & escape fraction & train Pearson per feature, last epoch (mean) \\
    \midrule
    \texttt{M\_zs\_lr3e-4\_mse} & 4 & 4 & \textbf{0.147 (0.025)} & 4 of 4 & 0.47 \\
    \texttt{M\_xattn\_lr3e-4} & 2 & 2 & 0.141 (0.010) & 2 of 2 & 0.71 \\
    \texttt{M\_pool\_lr3e-4} & 3 & 3 & 0.100 (0.072) & 2 of 3 & 0.39 \\
    \texttt{M\_zs\_lr3e-4} & 3 & 3 & 0.087 (0.084) & 3 of 3 & 0.35 \\
    \texttt{M\_cls\_graphkl1\_lr3e-4} & 1 & 1 & 0.077 & 0 of 1 & 0.12 \\
    \texttt{M\_cls\_graphkl1e-2\_lr3e-4} & 2 & 2 & 0.077 (0.031) & 2 of 2 & 0.49 \\
    \texttt{M\_zs\_crps\_lr1e-3} & 4 & 4 & 0.043 (0.018) & 0 of 4 & 0.10 \\
    \texttt{M\_cls\_crps\_lr1e-3} & 4 & 4 & 0.042 (0.016) & 0 of 4 & 0.10 \\
    \texttt{M\_cls\_lr3e-4} & 4 & 4 & 0.040 (0.016) & 0 of 4 & 0.09 \\
    \texttt{M\_cls\_lr1e-3\_mse} & 4 & 4 & 0.040 (0.018) & 0 of 4 & 0.10 \\
    \texttt{M\_zs\_lr1e-3\_mse} & 4 & 4 & 0.040 (0.016) & 0 of 4 & 0.10 \\
    \texttt{M\_cls\_lr3e-4\_mse} & 4 & 4 & 0.038 (0.017) & 0 of 4 & 0.07 \\
    \texttt{M\_cls\_lr3e-3} & 3 & 3 & 0.037 (0.004) & 0 of 3 & 0.11 \\
    \texttt{M\_zs\_lr3e-3} & 3 & 3 & 0.036 (0.005) & 0 of 3 & 0.11 \\
    \texttt{M\_pool\_lr3e-3} & 3 & 3 & 0.036 (0.004) & 0 of 3 & 0.11 \\
    \texttt{M\_xattn\_lr3e-3} & 3 & 3 & 0.035 (0.004) & 0 of 3 & 0.11 \\
    \texttt{M\_cls\_lr1e-3} & 3 & 3 & 0.035 (0.004) & 0 of 3 & 0.10 \\
    \texttt{M\_zs\_lr1e-3} & 3 & 3 & 0.035 (0.004) & 0 of 3 & 0.08 \\
    \texttt{M\_pool\_lr1e-3} & 3 & 3 & 0.034 (0.004) & 0 of 3 & 0.09 \\
    \texttt{M\_xattn\_lr1e-3} & 3 & 3 & 0.029 (0.004) & 0 of 3 & 0.07 \\
    \texttt{M\_zs\_lr1e-3\_pearson\_mse} & 4 & 4 & 0.000 (0.000) & 0 of 4 & 0.00 \\
    \texttt{M\_cls\_lr1e-3\_pearson\_mse} & 4 & 4 & 0.000 (0.000) & 0 of 4 & 0.00 \\
    \texttt{M\_pool\_b64\_lr3e-4} & 6 & 0 & partial & 0 of 6 & 0.11 \\
    \texttt{M\_zs\_b128\_lr3e-4} & 12 & 0 & partial & 3 of 12 & 0.24 \\
    \texttt{M\_zs\_b32\_lr3e-4} & 6 & 0 & partial & 0 of 6 & 0.11 \\
    \texttt{M\_pool\_b32\_lr3e-4} & 6 & 0 & partial & 0 of 6 & 0.13 \\
    \texttt{M\_pool\_b128\_lr3e-4} & 12 & 0 & partial & 0 of 12 & 0.09 \\
    \texttt{M\_zs\_b64\_lr3e-4} & 6 & 0 & partial & 1 of 6 & 0.23 \\
    \bottomrule
  \end{tabular}
\end{table}


\pfstep{The plateau and the two numbers that read it} Every run's validation score is read on
the window; whether a run left the marginal is read from the eval-mode training Pearson
crossing 0.3, which the plateau never approaches (0.04 to 0.12) and an escaped run passes
within about 50 epochs of leaving it. The two agree: no cell has a high window without
escapes, and the escaped runs of the plain \texttt{M\_cls} at lr $3\times10^{-4}$ (4 of 12,
over four splits and three initializations) score 0.13 to 0.18 where their siblings score
0.03. An escape is therefore a property of the run and not of the split, and a cell is
read as its escape fraction first and its window second.

\pfstep{Readout and loss interact} MSE on the CLS readout escaped 0 of 8 at two rates; the
pinball on the gene-token readout escaped 8 of 10 but at epochs 230 to 770, so that cell's
window ranges 0.04 to 0.21 by when the escape came; MSE on the gene-token readout escaped
12 of 12 by epoch 410 and gives the smallest spread of any cell. The cross-attention
readout escapes under the pinball (4 of 4, $0.121 \pm 0.029$) and not under MSE (1 of 4).
Hypothesis, untested: $z_S$ hands the head a vector whose scale the perturbation operator
sets directly, where the perturbed CLS must first learn to attend to its one key, and the
MSE gradient on the mean prediction is larger than the pinball's near the marginal, so it
moves the operator sooner; the cross-attention head has its own path to the tokens and
does not need either.

\pfstep{Learning rate} Twenty-seven cells at $10^{-3}$ or $3\times10^{-3}$, with and without
warmup or restarts, produced no escape. At $10^{-4}$ the plain CLS escaped 7 of 8 and the
gene-token readout with MSE 4 of 4 ($0.133 \pm 0.031$, below its $3\times10^{-4}$ reading);
the lower rate is a second route off the plateau for the CLS and not an improvement for
the cell that already leaves it. The 1{,}000-epoch runs peaked between epochs 400 and 600
and fell afterward (\cref{tab:morph-v25}, the json carries both windows), so the 500-epoch
budget stands.

\pfstep{Graph prior} The hard attention mask on layer 1 is the v24 default. Replacing it by
the KL regularizer at $10^{-2}$ made the CLS escape 4 of 4 at a window of 0.093; at 1 the
model escaped 1 of 4 and was still rising at 499. The prior is a third lever, weaker than the
readout; the stacked cell (MSE, KL $10^{-2}$, lr $10^{-3}$ on \texttt{M\_zs}) is the only cell
at lr $10^{-3}$ that escaped (4 of 4, 0.091), so the KL also widens the rate the model
tolerates.

\pfstep{Batch} On Delta, \texttt{M\_zs} under the pinball at batch 128 escaped 4 of 6 and at
batch 64 2 of 3, at batch 32 0 of 3; \texttt{M\_pool} escaped 0 of 12 at every batch, so the
two v24 \texttt{M\_pool} escapes at 1{,}200 epochs were the tail of a slow process and the
pool is not a reliable readout.

\pfstep{The contrastive term}
\label{sec:align}
\file{torchcell/losses/latent_align.py} adds, for every pair of supervised strains in a batch,
the squared gap between the cosine similarity of their latent vectors (the perturbed CLS,
$z_S$, or both; batch-centered) and the Pearson correlation of their label vectors over the
features both have measured, weighted $\lambda$ beside the data loss. The marginal
solution, every strain the same vector, scores a cosine of one for every pair, so the term
was expected to penalize the plateau from the first step. It did not: its own training value
stayed at 0.93 to 0.97 for 200 epochs and more in every run (the centered cosines of
near-identical latents are $\pm 1$), its validation value never fell, and the five morphology
cells ($\lambda$ 0.01, 0.1 and 1 on the CLS, one without centering; 0.1 and 1 on
$[\mathrm{CLS} \,\|\, z_S]$) sit at 0.039 to 0.047 with 0 of 20 escapes where the same
\texttt{M\_zs} without the term escaped 4 of 4. Hypothesis, untested: normalizing tiny
centered differences yields large gradients that the clip absorbs together with the data
gradient. On the knockout proteome (\cref{tab:prot-v26}: the full Messner store, 4{,}281
genotypes, the v21 small trunk at batch 128, 600 epochs, window epochs 500 to 599 declared
with the launch) the reference reads $0.062 \pm 0.010$, the level of v14 and v21, and the
term reads $-0.008$ (2 of 4 splits up) and $-0.010$ (1 of 4). Runs in
\cref{tab:wandb-runs-prot-v26}; the view is
\href{https://wandb.ai/zhao-group/torchcell_019_prot_v26?nw=9nsa6vtv42n}{torchcell\_019\_prot\_v26}.

%% GENERATED by experiments/019-simb-multimodal/scripts/prot_v26_readout.py. Do not edit by hand.
%% SOURCE: experiments/019-simb-multimodal/results/prot_v26_readout.json
\begin{table}[htbp]
  \centering
  \footnotesize
  \caption[Round v26, proteome alignment]{Round v26 on the knockout proteome (Messner 2023, every proteome genotype), genotype only, v21 small trunk at batch 128, 600 epochs. The window is the mean validation Pearson per feature over epochs 500 to 599 of finished runs; the paired difference is against the reference on the same split seed; the training fit is the eval-mode training Pearson at the last logged epoch. Bold, the best window mean. \src{experiments/019-simb-multimodal/scripts/prot_v26_readout.py}}
  \label{tab:prot-v26}
  \setlength{\tabcolsep}{3pt}
  \begin{tabular}{@{}>{\raggedright\arraybackslash}p{40mm}*{5}{>{\raggedleft\arraybackslash}p{26mm}}@{}}
    \toprule
    cell & runs (n) & finished (n) & window mean val Pearson per feature (s.d.) & paired difference vs reference, mean (splits up) & train Pearson per feature, last epoch (mean) \\
    \midrule
    \texttt{S\_prot\_ref} & 4 & 4 & \textbf{0.062 (0.010)} & +0.000 (0 of 4) & 0.51 \\
    \texttt{S\_prot\_align0.1\_srczs} & 4 & 4 & 0.055 (0.018) & -0.008 (2 of 4) & 0.50 \\
    \texttt{S\_prot\_align1\_srczs} & 4 & 4 & 0.053 (0.012) & -0.010 (1 of 4) & 0.40 \\
    \bottomrule
  \end{tabular}
\end{table}


\pfstep{A loss that was NaN from the first step} The eight \texttt{pearson\_mse} runs of
v25c logged NaN at every epoch. The per-feature Pearson loss took the square root of the
product of two column norms and clamped afterward; a global head at initialization emits
one vector for every strain (the perturbation operator's ReZero gate starts at zero), so a
prediction column's norm is exactly zero and the backward of the square root at zero is
infinite. The clamp now sits inside the root (\file{torchcell/losses/distributional.py}, with a
regression test on a mixed batch); the rerun trains and sits on the marginal ($0.040 \pm
0.016$, 0 of 8, at lr $10^{-3}$, the rate no cell escapes at, so the loss itself is not yet
tested at $3\times10^{-4}$).

\pfstep{What this changes for the figure} Panel c of the SI sketch (\cref{fig:si-mockup})
reads the genotype-only bar at 0.046; it now reads 0.149, within the seed spread of the
proteome-as-input arm (0.145) and above the sequence ridge. The claim that only a revealed
measurement moves morphology no longer holds as stated: the readout did, by three times.
Two consequences, both open. The measured-input arms of v24 (\texttt{CM\_prot},
\texttt{CM\_expr} and their permuted controls) were read off the CLS readout that never
escapes, so their gain over genotype only was measured against a model on the plateau; the
paired design has to be rerun on the \texttt{M\_zs} with MSE base before panel d's
differences mean what the caption says (proposal, not run). And the ceiling on these
features is 0.847, so 0.149 remains a sixth of what is reproducible; whether the escaped
model's errors are the same features the ridge misses is unmeasured.

%% GENERATED by experiments/019-simb-multimodal/scripts/morph_v25_readout.py. Do not edit by hand.
%% SOURCE: experiments/019-simb-multimodal/results/morph_v25_readout.json
\begin{table}[htbp]
  \centering
  \footnotesize
  \caption[W\&B runs, morph-v25]{Every run of rounds v25 and v26 (project torchcell\_019\_morph\_v25) by cell, linked to its W\&B page. \src{experiments/019-simb-multimodal/scripts/morph_v25_readout.py}}
  \label{tab:wandb-runs-morph-v25}
  \begin{tabular}{@{}>{\raggedright\arraybackslash}p{60mm}>{\raggedright\arraybackslash}p{118mm}@{}}
    \toprule
    cell & runs (split seed, init seed: W\&B run id) \\
    \midrule
    \texttt{M\_cls\_align0.01\_lr3e-4\_srccls} & s0i0: \href{https://wandb.ai/zhao-group/torchcell_019_morph_v25/runs/u3x6lga4}{\texttt{u3x6lga4}}, s1i0: \href{https://wandb.ai/zhao-group/torchcell_019_morph_v25/runs/y6pbwxo9}{\texttt{y6pbwxo9}}, s2i0: \href{https://wandb.ai/zhao-group/torchcell_019_morph_v25/runs/zcjri9ds}{\texttt{zcjri9ds}}, s3i0: \href{https://wandb.ai/zhao-group/torchcell_019_morph_v25/runs/j7dggjnt}{\texttt{j7dggjnt}} \\
    \texttt{M\_cls\_align0.1\_lr3e-4\_srccls} & s0i0: \href{https://wandb.ai/zhao-group/torchcell_019_morph_v25/runs/eq5nqrah}{\texttt{eq5nqrah}}, s1i0: \href{https://wandb.ai/zhao-group/torchcell_019_morph_v25/runs/egc1c1s0}{\texttt{egc1c1s0}}, s2i0: \href{https://wandb.ai/zhao-group/torchcell_019_morph_v25/runs/kcttphqs}{\texttt{kcttphqs}}, s3i0: \href{https://wandb.ai/zhao-group/torchcell_019_morph_v25/runs/i6pxf72w}{\texttt{i6pxf72w}} \\
    \texttt{M\_cls\_align1\_lr3e-4\_nocenter\_srccls} & s0i0: \href{https://wandb.ai/zhao-group/torchcell_019_morph_v25/runs/s8w6bvep}{\texttt{s8w6bvep}}, s1i0: \href{https://wandb.ai/zhao-group/torchcell_019_morph_v25/runs/xj2615ek}{\texttt{xj2615ek}}, s2i0: \href{https://wandb.ai/zhao-group/torchcell_019_morph_v25/runs/d4k9vm9f}{\texttt{d4k9vm9f}}, s3i0: \href{https://wandb.ai/zhao-group/torchcell_019_morph_v25/runs/cqdvx2uo}{\texttt{cqdvx2uo}} \\
    \texttt{M\_cls\_align1\_lr3e-4\_srccls} & s0i0: \href{https://wandb.ai/zhao-group/torchcell_019_morph_v25/runs/g433p1xq}{\texttt{g433p1xq}}, s1i0: \href{https://wandb.ai/zhao-group/torchcell_019_morph_v25/runs/ssxz0zwx}{\texttt{ssxz0zwx}}, s2i0: \href{https://wandb.ai/zhao-group/torchcell_019_morph_v25/runs/57ejlkdb}{\texttt{57ejlkdb}}, s3i0: \href{https://wandb.ai/zhao-group/torchcell_019_morph_v25/runs/2dcs65s8}{\texttt{2dcs65s8}} \\
    \texttt{M\_cls\_crps\_lr1e-3} & s0i0: \href{https://wandb.ai/zhao-group/torchcell_019_morph_v25/runs/p791cvew}{\texttt{p791cvew}}, s1i0: \href{https://wandb.ai/zhao-group/torchcell_019_morph_v25/runs/fwuh3jqd}{\texttt{fwuh3jqd}}, s2i0: \href{https://wandb.ai/zhao-group/torchcell_019_morph_v25/runs/ts8nojhj}{\texttt{ts8nojhj}}, s3i0: \href{https://wandb.ai/zhao-group/torchcell_019_morph_v25/runs/ksgu9h8n}{\texttt{ksgu9h8n}} \\
    \texttt{M\_cls\_graphkl1\_lr3e-4} & s0i0: \href{https://wandb.ai/zhao-group/torchcell_019_morph_v25/runs/z7u8yavf}{\texttt{z7u8yavf}}, s1i0: \href{https://wandb.ai/zhao-group/torchcell_019_morph_v25/runs/2qx3vvjl}{\texttt{2qx3vvjl}}, s2i0: \href{https://wandb.ai/zhao-group/torchcell_019_morph_v25/runs/zgbrbfaz}{\texttt{zgbrbfaz}}, s3i0: \href{https://wandb.ai/zhao-group/torchcell_019_morph_v25/runs/f4dbf383}{\texttt{f4dbf383}} \\
    \texttt{M\_cls\_graphkl1e-2\_lr3e-4} & s0i0: \href{https://wandb.ai/zhao-group/torchcell_019_morph_v25/runs/v6z10hp2}{\texttt{v6z10hp2}}, s1i0: \href{https://wandb.ai/zhao-group/torchcell_019_morph_v25/runs/2va89u0l}{\texttt{2va89u0l}}, s2i0: \href{https://wandb.ai/zhao-group/torchcell_019_morph_v25/runs/7k72hf52}{\texttt{7k72hf52}}, s3i0: \href{https://wandb.ai/zhao-group/torchcell_019_morph_v25/runs/nfas55te}{\texttt{nfas55te}} \\
    \texttt{M\_cls\_lr1e-3} & s0i0: \href{https://wandb.ai/zhao-group/torchcell_019_morph_v25/runs/yvqzu67v}{\texttt{yvqzu67v}}, s1i0: \href{https://wandb.ai/zhao-group/torchcell_019_morph_v25/runs/kjynf7ag}{\texttt{kjynf7ag}}, s2i0: \href{https://wandb.ai/zhao-group/torchcell_019_morph_v25/runs/va0buthp}{\texttt{va0buthp}}, s3i0: \href{https://wandb.ai/zhao-group/torchcell_019_morph_v25/runs/p3igpx7j}{\texttt{p3igpx7j}} \\
    \texttt{M\_cls\_lr1e-3\_mse} & s0i0: \href{https://wandb.ai/zhao-group/torchcell_019_morph_v25/runs/ewqbxwh5}{\texttt{ewqbxwh5}}, s1i0: \href{https://wandb.ai/zhao-group/torchcell_019_morph_v25/runs/uzehb488}{\texttt{uzehb488}}, s2i0: \href{https://wandb.ai/zhao-group/torchcell_019_morph_v25/runs/6cfm75im}{\texttt{6cfm75im}}, s3i0: \href{https://wandb.ai/zhao-group/torchcell_019_morph_v25/runs/8dy1hufk}{\texttt{8dy1hufk}} \\
    \texttt{M\_cls\_lr1e-3\_pearson\_mse} & s0i0: \href{https://wandb.ai/zhao-group/torchcell_019_morph_v25/runs/c2o3nroj}{\texttt{c2o3nroj}}, s0i0: \href{https://wandb.ai/zhao-group/torchcell_019_morph_v25/runs/8ny19in9}{\texttt{8ny19in9}}, s1i0: \href{https://wandb.ai/zhao-group/torchcell_019_morph_v25/runs/keftxjac}{\texttt{keftxjac}}, s1i0: \href{https://wandb.ai/zhao-group/torchcell_019_morph_v25/runs/0cam21uu}{\texttt{0cam21uu}}, s2i0: \href{https://wandb.ai/zhao-group/torchcell_019_morph_v25/runs/la5y087j}{\texttt{la5y087j}}, s2i0: \href{https://wandb.ai/zhao-group/torchcell_019_morph_v25/runs/xy45v4b1}{\texttt{xy45v4b1}}, s3i0: \href{https://wandb.ai/zhao-group/torchcell_019_morph_v25/runs/et9tubov}{\texttt{et9tubov}}, s3i0: \href{https://wandb.ai/zhao-group/torchcell_019_morph_v25/runs/qr2cih8n}{\texttt{qr2cih8n}} \\
    \texttt{M\_cls\_lr1e-3\_restarts300} & s0i0: \href{https://wandb.ai/zhao-group/torchcell_019_morph_v25/runs/2f06t5z2}{\texttt{2f06t5z2}}, s1i0: \href{https://wandb.ai/zhao-group/torchcell_019_morph_v25/runs/hsn856u0}{\texttt{hsn856u0}}, s2i0: \href{https://wandb.ai/zhao-group/torchcell_019_morph_v25/runs/vz5eh0ew}{\texttt{vz5eh0ew}}, s3i0: \href{https://wandb.ai/zhao-group/torchcell_019_morph_v25/runs/ltzfvsu9}{\texttt{ltzfvsu9}} \\
    \texttt{M\_cls\_lr1e-4} & s0i0: \href{https://wandb.ai/zhao-group/torchcell_019_morph_v25/runs/b3b1kzqm}{\texttt{b3b1kzqm}}, s0i0: \href{https://wandb.ai/zhao-group/torchcell_019_morph_v25/runs/bv7fj78p}{\texttt{bv7fj78p}}, s1i0: \href{https://wandb.ai/zhao-group/torchcell_019_morph_v25/runs/1wqug3sz}{\texttt{1wqug3sz}}, s1i0: \href{https://wandb.ai/zhao-group/torchcell_019_morph_v25/runs/sgyrzswe}{\texttt{sgyrzswe}}, s2i0: \href{https://wandb.ai/zhao-group/torchcell_019_morph_v25/runs/kd1ko3tt}{\texttt{kd1ko3tt}}, s2i0: \href{https://wandb.ai/zhao-group/torchcell_019_morph_v25/runs/xhle3dsh}{\texttt{xhle3dsh}}, s3i0: \href{https://wandb.ai/zhao-group/torchcell_019_morph_v25/runs/f317sr5g}{\texttt{f317sr5g}}, s3i0: \href{https://wandb.ai/zhao-group/torchcell_019_morph_v25/runs/qswl42s2}{\texttt{qswl42s2}} \\
    \texttt{M\_cls\_lr1e-4\_mse} & s0i0: \href{https://wandb.ai/zhao-group/torchcell_019_morph_v25/runs/1y19vm52}{\texttt{1y19vm52}}, s1i0: \href{https://wandb.ai/zhao-group/torchcell_019_morph_v25/runs/p75djmj1}{\texttt{p75djmj1}}, s2i0: \href{https://wandb.ai/zhao-group/torchcell_019_morph_v25/runs/ql5mozzz}{\texttt{ql5mozzz}}, s3i0: \href{https://wandb.ai/zhao-group/torchcell_019_morph_v25/runs/0yfimmsc}{\texttt{0yfimmsc}} \\
    \texttt{M\_cls\_lr3e-3} & s0i0: \href{https://wandb.ai/zhao-group/torchcell_019_morph_v25/runs/hrxu1n0d}{\texttt{hrxu1n0d}}, s1i0: \href{https://wandb.ai/zhao-group/torchcell_019_morph_v25/runs/24w4aajx}{\texttt{24w4aajx}}, s2i0: \href{https://wandb.ai/zhao-group/torchcell_019_morph_v25/runs/7dyr9g50}{\texttt{7dyr9g50}} \\
    \texttt{M\_cls\_lr3e-4} & s0i0: \href{https://wandb.ai/zhao-group/torchcell_019_morph_v25/runs/d5n9hfc0}{\texttt{d5n9hfc0}}, s0i1: \href{https://wandb.ai/zhao-group/torchcell_019_morph_v25/runs/ehcts6qf}{\texttt{ehcts6qf}}, s0i2: \href{https://wandb.ai/zhao-group/torchcell_019_morph_v25/runs/ppye6rr9}{\texttt{ppye6rr9}}, s1i0: \href{https://wandb.ai/zhao-group/torchcell_019_morph_v25/runs/rqntj7qs}{\texttt{rqntj7qs}}, s1i1: \href{https://wandb.ai/zhao-group/torchcell_019_morph_v25/runs/frrow7ej}{\texttt{frrow7ej}}, s1i2: \href{https://wandb.ai/zhao-group/torchcell_019_morph_v25/runs/x1hekr4h}{\texttt{x1hekr4h}}, s2i0: \href{https://wandb.ai/zhao-group/torchcell_019_morph_v25/runs/g62elvtb}{\texttt{g62elvtb}}, s2i1: \href{https://wandb.ai/zhao-group/torchcell_019_morph_v25/runs/y7m8y86j}{\texttt{y7m8y86j}}, s2i2: \href{https://wandb.ai/zhao-group/torchcell_019_morph_v25/runs/rgo09hlx}{\texttt{rgo09hlx}}, s3i0: \href{https://wandb.ai/zhao-group/torchcell_019_morph_v25/runs/ycvvxa07}{\texttt{ycvvxa07}}, s3i1: \href{https://wandb.ai/zhao-group/torchcell_019_morph_v25/runs/bzp6wcqc}{\texttt{bzp6wcqc}}, s3i2: \href{https://wandb.ai/zhao-group/torchcell_019_morph_v25/runs/vyskzxbg}{\texttt{vyskzxbg}} \\
    \texttt{M\_cls\_lr3e-4\_mse} & s0i0: \href{https://wandb.ai/zhao-group/torchcell_019_morph_v25/runs/613ww08p}{\texttt{613ww08p}}, s1i0: \href{https://wandb.ai/zhao-group/torchcell_019_morph_v25/runs/6kab7to5}{\texttt{6kab7to5}}, s2i0: \href{https://wandb.ai/zhao-group/torchcell_019_morph_v25/runs/69p5zpql}{\texttt{69p5zpql}}, s3i0: \href{https://wandb.ai/zhao-group/torchcell_019_morph_v25/runs/286cj131}{\texttt{286cj131}} \\
    \texttt{M\_pool\_b128\_lr3e-4} & s0i1: \href{https://wandb.ai/zhao-group/torchcell_019_morph_v25/runs/y08llh8a}{\texttt{y08llh8a}}, s0i2: \href{https://wandb.ai/zhao-group/torchcell_019_morph_v25/runs/4iode6i5}{\texttt{4iode6i5}}, s1i1: \href{https://wandb.ai/zhao-group/torchcell_019_morph_v25/runs/0gk17xzn}{\texttt{0gk17xzn}}, s1i2: \href{https://wandb.ai/zhao-group/torchcell_019_morph_v25/runs/4kdi82mu}{\texttt{4kdi82mu}}, s2i1: \href{https://wandb.ai/zhao-group/torchcell_019_morph_v25/runs/q1t1lrrd}{\texttt{q1t1lrrd}}, s2i2: \href{https://wandb.ai/zhao-group/torchcell_019_morph_v25/runs/up80kss7}{\texttt{up80kss7}} \\
    \texttt{M\_pool\_b32\_lr3e-4} & s0i0: \href{https://wandb.ai/zhao-group/torchcell_019_morph_v25/runs/u2oy5100}{\texttt{u2oy5100}}, s1i0: \href{https://wandb.ai/zhao-group/torchcell_019_morph_v25/runs/0ifuihme}{\texttt{0ifuihme}}, s2i0: \href{https://wandb.ai/zhao-group/torchcell_019_morph_v25/runs/f03vxtl0}{\texttt{f03vxtl0}} \\
    \texttt{M\_pool\_b64\_lr3e-4} & s0i0: \href{https://wandb.ai/zhao-group/torchcell_019_morph_v25/runs/g17i75w9}{\texttt{g17i75w9}}, s1i0: \href{https://wandb.ai/zhao-group/torchcell_019_morph_v25/runs/wygenmjx}{\texttt{wygenmjx}}, s2i0: \href{https://wandb.ai/zhao-group/torchcell_019_morph_v25/runs/uafbcuak}{\texttt{uafbcuak}} \\
    \texttt{M\_pool\_lr1e-3} & s0i0: \href{https://wandb.ai/zhao-group/torchcell_019_morph_v25/runs/aubga888}{\texttt{aubga888}}, s1i0: \href{https://wandb.ai/zhao-group/torchcell_019_morph_v25/runs/zo9t60bj}{\texttt{zo9t60bj}}, s2i0: \href{https://wandb.ai/zhao-group/torchcell_019_morph_v25/runs/8vwr6zi5}{\texttt{8vwr6zi5}}, s3i0: \href{https://wandb.ai/zhao-group/torchcell_019_morph_v25/runs/xoed8q6h}{\texttt{xoed8q6h}} \\
    \texttt{M\_pool\_lr3e-3} & s0i0: \href{https://wandb.ai/zhao-group/torchcell_019_morph_v25/runs/3ho7pt2i}{\texttt{3ho7pt2i}}, s1i0: \href{https://wandb.ai/zhao-group/torchcell_019_morph_v25/runs/p490gtjh}{\texttt{p490gtjh}}, s2i0: \href{https://wandb.ai/zhao-group/torchcell_019_morph_v25/runs/rrcjlkkg}{\texttt{rrcjlkkg}} \\
    \texttt{M\_pool\_lr3e-4} & s0i0: \href{https://wandb.ai/zhao-group/torchcell_019_morph_v25/runs/sxketfkz}{\texttt{sxketfkz}}, s1i0: \href{https://wandb.ai/zhao-group/torchcell_019_morph_v25/runs/rnmy125d}{\texttt{rnmy125d}}, s2i0: \href{https://wandb.ai/zhao-group/torchcell_019_morph_v25/runs/mhbcyyz1}{\texttt{mhbcyyz1}}, s3i0: \href{https://wandb.ai/zhao-group/torchcell_019_morph_v25/runs/ddktblke}{\texttt{ddktblke}} \\
    \texttt{M\_xattn\_lr1e-3} & s0i0: \href{https://wandb.ai/zhao-group/torchcell_019_morph_v25/runs/lsyq6mv2}{\texttt{lsyq6mv2}}, s1i0: \href{https://wandb.ai/zhao-group/torchcell_019_morph_v25/runs/zhljnci5}{\texttt{zhljnci5}}, s2i0: \href{https://wandb.ai/zhao-group/torchcell_019_morph_v25/runs/pex46wxp}{\texttt{pex46wxp}} \\
    \texttt{M\_xattn\_lr3e-3} & s0i0: \href{https://wandb.ai/zhao-group/torchcell_019_morph_v25/runs/6iehpw8a}{\texttt{6iehpw8a}}, s1i0: \href{https://wandb.ai/zhao-group/torchcell_019_morph_v25/runs/nlwpx250}{\texttt{nlwpx250}}, s2i0: \href{https://wandb.ai/zhao-group/torchcell_019_morph_v25/runs/wk1yju9m}{\texttt{wk1yju9m}} \\
    \texttt{M\_xattn\_lr3e-4} & s0i0: \href{https://wandb.ai/zhao-group/torchcell_019_morph_v25/runs/790bz2gu}{\texttt{790bz2gu}}, s1i0: \href{https://wandb.ai/zhao-group/torchcell_019_morph_v25/runs/0p2hkd3l}{\texttt{0p2hkd3l}}, s2i0: \href{https://wandb.ai/zhao-group/torchcell_019_morph_v25/runs/zd2h4qzn}{\texttt{zd2h4qzn}}, s3i0: \href{https://wandb.ai/zhao-group/torchcell_019_morph_v25/runs/qxcaf37i}{\texttt{qxcaf37i}} \\
    \texttt{M\_xattn\_lr3e-4\_mse} & s0i0: \href{https://wandb.ai/zhao-group/torchcell_019_morph_v25/runs/q7w8mc76}{\texttt{q7w8mc76}}, s1i0: \href{https://wandb.ai/zhao-group/torchcell_019_morph_v25/runs/5w8c9dnp}{\texttt{5w8c9dnp}}, s2i0: \href{https://wandb.ai/zhao-group/torchcell_019_morph_v25/runs/ui4s3amu}{\texttt{ui4s3amu}}, s3i0: \href{https://wandb.ai/zhao-group/torchcell_019_morph_v25/runs/l1fcrih4}{\texttt{l1fcrih4}} \\
    \texttt{M\_zs\_align0.1\_lr3e-4\_mse\_srccls\_zs} & s0i0: \href{https://wandb.ai/zhao-group/torchcell_019_morph_v25/runs/cbvejneq}{\texttt{cbvejneq}} \\
    \texttt{M\_zs\_align0.1\_lr3e-4\_srccls\_zs} & s0i0: \href{https://wandb.ai/zhao-group/torchcell_019_morph_v25/runs/tw2jjokw}{\texttt{tw2jjokw}}, s1i0: \href{https://wandb.ai/zhao-group/torchcell_019_morph_v25/runs/4fvf9kaz}{\texttt{4fvf9kaz}}, s2i0: \href{https://wandb.ai/zhao-group/torchcell_019_morph_v25/runs/n9417cwz}{\texttt{n9417cwz}}, s3i0: \href{https://wandb.ai/zhao-group/torchcell_019_morph_v25/runs/g72pfvgi}{\texttt{g72pfvgi}} \\
    \texttt{M\_zs\_align1\_lr3e-4\_srccls\_zs} & s0i0: \href{https://wandb.ai/zhao-group/torchcell_019_morph_v25/runs/zflos0pi}{\texttt{zflos0pi}}, s1i0: \href{https://wandb.ai/zhao-group/torchcell_019_morph_v25/runs/qsgazw26}{\texttt{qsgazw26}}, s2i0: \href{https://wandb.ai/zhao-group/torchcell_019_morph_v25/runs/dw7ctebl}{\texttt{dw7ctebl}}, s3i0: \href{https://wandb.ai/zhao-group/torchcell_019_morph_v25/runs/o6afkykc}{\texttt{o6afkykc}} \\
    \texttt{M\_zs\_b128\_lr3e-4} & s0i1: \href{https://wandb.ai/zhao-group/torchcell_019_morph_v25/runs/xcboejey}{\texttt{xcboejey}}, s0i2: \href{https://wandb.ai/zhao-group/torchcell_019_morph_v25/runs/aassz7l8}{\texttt{aassz7l8}}, s1i1: \href{https://wandb.ai/zhao-group/torchcell_019_morph_v25/runs/pkuk57yd}{\texttt{pkuk57yd}}, s1i2: \href{https://wandb.ai/zhao-group/torchcell_019_morph_v25/runs/04p8c0q5}{\texttt{04p8c0q5}}, s2i1: \href{https://wandb.ai/zhao-group/torchcell_019_morph_v25/runs/p65o7udc}{\texttt{p65o7udc}}, s2i2: \href{https://wandb.ai/zhao-group/torchcell_019_morph_v25/runs/ex6xkhwk}{\texttt{ex6xkhwk}} \\
    \texttt{M\_zs\_b32\_lr3e-4} & s0i0: \href{https://wandb.ai/zhao-group/torchcell_019_morph_v25/runs/o97mckca}{\texttt{o97mckca}}, s1i0: \href{https://wandb.ai/zhao-group/torchcell_019_morph_v25/runs/na1q1jwb}{\texttt{na1q1jwb}}, s2i0: \href{https://wandb.ai/zhao-group/torchcell_019_morph_v25/runs/47bwbtnx}{\texttt{47bwbtnx}} \\
    \texttt{M\_zs\_b64\_lr3e-4} & s0i0: \href{https://wandb.ai/zhao-group/torchcell_019_morph_v25/runs/uk6olozo}{\texttt{uk6olozo}}, s1i0: \href{https://wandb.ai/zhao-group/torchcell_019_morph_v25/runs/6varwr29}{\texttt{6varwr29}}, s2i0: \href{https://wandb.ai/zhao-group/torchcell_019_morph_v25/runs/cc32b3h3}{\texttt{cc32b3h3}} \\
    \texttt{M\_zs\_crps\_lr1e-3} & s0i0: \href{https://wandb.ai/zhao-group/torchcell_019_morph_v25/runs/o3s3f2xj}{\texttt{o3s3f2xj}}, s1i0: \href{https://wandb.ai/zhao-group/torchcell_019_morph_v25/runs/x5rhb6ng}{\texttt{x5rhb6ng}}, s2i0: \href{https://wandb.ai/zhao-group/torchcell_019_morph_v25/runs/qv7t9i99}{\texttt{qv7t9i99}}, s3i0: \href{https://wandb.ai/zhao-group/torchcell_019_morph_v25/runs/cwzt3f9i}{\texttt{cwzt3f9i}} \\
    \texttt{M\_zs\_graphkl1e-2\_lr1e-3\_mse\_stack} & s0i0: \href{https://wandb.ai/zhao-group/torchcell_019_morph_v25/runs/y5uejwn7}{\texttt{y5uejwn7}}, s1i0: \href{https://wandb.ai/zhao-group/torchcell_019_morph_v25/runs/f5mi7n3f}{\texttt{f5mi7n3f}}, s2i0: \href{https://wandb.ai/zhao-group/torchcell_019_morph_v25/runs/wcinelfs}{\texttt{wcinelfs}}, s3i0: \href{https://wandb.ai/zhao-group/torchcell_019_morph_v25/runs/bfpoupdj}{\texttt{bfpoupdj}} \\
    \texttt{M\_zs\_lr1e-3} & s0i0: \href{https://wandb.ai/zhao-group/torchcell_019_morph_v25/runs/e2v1y0ap}{\texttt{e2v1y0ap}}, s1i0: \href{https://wandb.ai/zhao-group/torchcell_019_morph_v25/runs/2osvaj68}{\texttt{2osvaj68}}, s2i0: \href{https://wandb.ai/zhao-group/torchcell_019_morph_v25/runs/nbw8ien1}{\texttt{nbw8ien1}} \\
    \texttt{M\_zs\_lr1e-3\_mse} & s0i0: \href{https://wandb.ai/zhao-group/torchcell_019_morph_v25/runs/vclpccrc}{\texttt{vclpccrc}}, s1i0: \href{https://wandb.ai/zhao-group/torchcell_019_morph_v25/runs/ufl3wbvb}{\texttt{ufl3wbvb}}, s2i0: \href{https://wandb.ai/zhao-group/torchcell_019_morph_v25/runs/w6uah4hd}{\texttt{w6uah4hd}}, s3i0: \href{https://wandb.ai/zhao-group/torchcell_019_morph_v25/runs/44y7yqus}{\texttt{44y7yqus}} \\
    \texttt{M\_zs\_lr1e-3\_pearson\_mse} & s0i0: \href{https://wandb.ai/zhao-group/torchcell_019_morph_v25/runs/lqgbvss8}{\texttt{lqgbvss8}}, s0i0: \href{https://wandb.ai/zhao-group/torchcell_019_morph_v25/runs/jk5035z4}{\texttt{jk5035z4}}, s1i0: \href{https://wandb.ai/zhao-group/torchcell_019_morph_v25/runs/tx33fwi1}{\texttt{tx33fwi1}}, s1i0: \href{https://wandb.ai/zhao-group/torchcell_019_morph_v25/runs/mwk8jx0e}{\texttt{mwk8jx0e}}, s2i0: \href{https://wandb.ai/zhao-group/torchcell_019_morph_v25/runs/bxp8o1bb}{\texttt{bxp8o1bb}}, s2i0: \href{https://wandb.ai/zhao-group/torchcell_019_morph_v25/runs/e2ph5a5z}{\texttt{e2ph5a5z}}, s3i0: \href{https://wandb.ai/zhao-group/torchcell_019_morph_v25/runs/uaadslgw}{\texttt{uaadslgw}}, s3i0: \href{https://wandb.ai/zhao-group/torchcell_019_morph_v25/runs/s2pk8zcs}{\texttt{s2pk8zcs}} \\
    \texttt{M\_zs\_lr1e-4\_mse} & s0i0: \href{https://wandb.ai/zhao-group/torchcell_019_morph_v25/runs/hd26qp3c}{\texttt{hd26qp3c}}, s1i0: \href{https://wandb.ai/zhao-group/torchcell_019_morph_v25/runs/yp8jul6f}{\texttt{yp8jul6f}}, s2i0: \href{https://wandb.ai/zhao-group/torchcell_019_morph_v25/runs/vhzc9o0z}{\texttt{vhzc9o0z}}, s3i0: \href{https://wandb.ai/zhao-group/torchcell_019_morph_v25/runs/453iatob}{\texttt{453iatob}} \\
    \texttt{M\_zs\_lr3e-3} & s0i0: \href{https://wandb.ai/zhao-group/torchcell_019_morph_v25/runs/wkxhqp4t}{\texttt{wkxhqp4t}}, s1i0: \href{https://wandb.ai/zhao-group/torchcell_019_morph_v25/runs/ii5kuym1}{\texttt{ii5kuym1}}, s2i0: \href{https://wandb.ai/zhao-group/torchcell_019_morph_v25/runs/f8hui5v0}{\texttt{f8hui5v0}} \\
    \texttt{M\_zs\_lr3e-4} & s0i0: \href{https://wandb.ai/zhao-group/torchcell_019_morph_v25/runs/9uf782za}{\texttt{9uf782za}}, s1i0: \href{https://wandb.ai/zhao-group/torchcell_019_morph_v25/runs/vwxswdcx}{\texttt{vwxswdcx}}, s2i0: \href{https://wandb.ai/zhao-group/torchcell_019_morph_v25/runs/dndezjdw}{\texttt{dndezjdw}}, s3i0: \href{https://wandb.ai/zhao-group/torchcell_019_morph_v25/runs/gpxcyqpp}{\texttt{gpxcyqpp}} \\
    \texttt{M\_zs\_lr3e-4\_mse} & s0i0: \href{https://wandb.ai/zhao-group/torchcell_019_morph_v25/runs/e9l0eg4d}{\texttt{e9l0eg4d}}, s0i1: \href{https://wandb.ai/zhao-group/torchcell_019_morph_v25/runs/i30q51ho}{\texttt{i30q51ho}}, s0i2: \href{https://wandb.ai/zhao-group/torchcell_019_morph_v25/runs/ee09tqow}{\texttt{ee09tqow}}, s1i0: \href{https://wandb.ai/zhao-group/torchcell_019_morph_v25/runs/vexd9lw6}{\texttt{vexd9lw6}}, s1i1: \href{https://wandb.ai/zhao-group/torchcell_019_morph_v25/runs/7dcm1auh}{\texttt{7dcm1auh}}, s1i2: \href{https://wandb.ai/zhao-group/torchcell_019_morph_v25/runs/7js4wqxr}{\texttt{7js4wqxr}}, s2i0: \href{https://wandb.ai/zhao-group/torchcell_019_morph_v25/runs/xc2h0i07}{\texttt{xc2h0i07}}, s2i1: \href{https://wandb.ai/zhao-group/torchcell_019_morph_v25/runs/cg129w4t}{\texttt{cg129w4t}}, s2i2: \href{https://wandb.ai/zhao-group/torchcell_019_morph_v25/runs/agx0bqk8}{\texttt{agx0bqk8}}, s3i0: \href{https://wandb.ai/zhao-group/torchcell_019_morph_v25/runs/amyu3c6c}{\texttt{amyu3c6c}}, s3i1: \href{https://wandb.ai/zhao-group/torchcell_019_morph_v25/runs/jts7zqsd}{\texttt{jts7zqsd}}, s3i2: \href{https://wandb.ai/zhao-group/torchcell_019_morph_v25/runs/iazgk2f1}{\texttt{iazgk2f1}} \\
    \bottomrule
  \end{tabular}
\end{table}

%% GENERATED by experiments/019-simb-multimodal/scripts/prot_v26_readout.py. Do not edit by hand.
%% SOURCE: experiments/019-simb-multimodal/results/prot_v26_readout.json
\begin{table}[htbp]
  \centering
  \footnotesize
  \caption[W\&B runs, prot-v26]{Every run of round prot\_v26 (project torchcell\_019\_prot\_v26) by cell, linked to its W\&B page. \src{experiments/019-simb-multimodal/scripts/prot_v26_readout.py}}
  \label{tab:wandb-runs-prot-v26}
  \begin{tabular}{@{}>{\raggedright\arraybackslash}p{60mm}>{\raggedright\arraybackslash}p{118mm}@{}}
    \toprule
    cell & runs (split seed, init seed: W\&B run id) \\
    \midrule
    \texttt{S\_prot\_align0.1\_srczs} & s0i0: \href{https://wandb.ai/zhao-group/torchcell_019_prot_v26/runs/k6yj2rge}{\texttt{k6yj2rge}}, s1i0: \href{https://wandb.ai/zhao-group/torchcell_019_prot_v26/runs/ox4dz2my}{\texttt{ox4dz2my}}, s2i0: \href{https://wandb.ai/zhao-group/torchcell_019_prot_v26/runs/zgauc2i7}{\texttt{zgauc2i7}}, s3i0: \href{https://wandb.ai/zhao-group/torchcell_019_prot_v26/runs/n9e07q46}{\texttt{n9e07q46}} \\
    \texttt{S\_prot\_align1\_srczs} & s0i0: \href{https://wandb.ai/zhao-group/torchcell_019_prot_v26/runs/prxu3tr4}{\texttt{prxu3tr4}}, s1i0: \href{https://wandb.ai/zhao-group/torchcell_019_prot_v26/runs/jfrdzy98}{\texttt{jfrdzy98}}, s2i0: \href{https://wandb.ai/zhao-group/torchcell_019_prot_v26/runs/5006fcs7}{\texttt{5006fcs7}}, s3i0: \href{https://wandb.ai/zhao-group/torchcell_019_prot_v26/runs/rf1xqiqn}{\texttt{rf1xqiqn}} \\
    \texttt{S\_prot\_ref} & s0i0: \href{https://wandb.ai/zhao-group/torchcell_019_prot_v26/runs/3lcj4hhr}{\texttt{3lcj4hhr}}, s1i0: \href{https://wandb.ai/zhao-group/torchcell_019_prot_v26/runs/meu8dlz1}{\texttt{meu8dlz1}}, s2i0: \href{https://wandb.ai/zhao-group/torchcell_019_prot_v26/runs/0ij48gqo}{\texttt{0ij48gqo}}, s3i0: \href{https://wandb.ai/zhao-group/torchcell_019_prot_v26/runs/lvdm0r08}{\texttt{lvdm0r08}} \\
    \bottomrule
  \end{tabular}
\end{table}

